% GENERATED FILE. Edit canonical lessons, then run npm run generate:books.
% canonical-source-hash: fnv1a64:7b5b46e7474add3f
% canonical-lessons: MR-C368-palan-karne, MR-C368-ullanghan-karne, MR-C368-bajavne, MR-C368-sopavne, MR-C368-hatalne

\chapter{Duties and Rules}
\label{ch:mr-duties-and-rules}
\hlchaptermodality{368}

\begin{chapteropening}
\textbf{By the end of this chapter:} \emph{I can talk about following rules, breaking them and doing one's duty.}

Complete the last lesson of chapter 368: I can talk about following rules, breaking them and doing one's duty.
\end{chapteropening}

\section[\texorpdfstring{\mr{पालन} \mr{करणे} (pālan karṇe)}{पालन करणे (pālan karṇe)}]{\mr{पालन} \mr{करणे} (pālan karṇe) — to comply with; to observe}

\label{lesson:MR-C368-palan-karne}



\begin{quote}
Before the new one: say the Marathi for a tendency, then the Marathi for social.
\end{quote}

\subsection*{You'll want to know: \mr{पालन} \mr{करणे}}
\textbf{\mr{पालन} \mr{करणे}} — \emph{pālan karṇe} — \textquotedblleft{}to comply with; to observe\textquotedblright{}.

\begin{quote}
\mr{सगळ्यांनी} \mr{वाहतुकीच्या} \mr{नियमांचं} \mr{पालन} \mr{केलं} \mr{पाहिजे}.
\end{quote}

\begin{quote}
\emph{sagaḷyānnī vāhtukīcyā niymāñcaṁ pālan kelaṁ pāhije.}
\end{quote}

\begin{quote}
Everyone must observe the traffic rules.
\end{quote}

Past, with -\mr{ने} (he …): \mr{त्याने} \mr{पालन} \mr{केलं}, \emph{tyāne pālan kelaṁ}.

The rule takes -\mr{चं}: \mr{नियमांचं} \mr{पालन} \mr{करणे}; also to bring up: \mr{मुलांचं} \mr{पालन}.

\begin{cousinweb}
\textbf{\mr{पालन} \mr{करणे}} (\emph{pālan karṇe}) joins the tatsama \mr{पालन}, Sanskrit \emph{pālana}, \textquotedblleft{}protecting, keeping\textquotedblright{} (CDIAL 8128), to \mr{करणे}, \textquotedblleft{}to do\textquotedblright{}.
\end{cousinweb}

\begin{grammarlens}[title={What you've built}]
One more word for giving reasons and opinions.
\end{grammarlens}

\subsection*{Your turn}
\begin{itemize}
\raggedright
  \item \emph{Say it:} \emph{pālan karṇe}
  \item \emph{Say it:} \emph{pālan karṇe}, once more
  \item \emph{Say it:} \emph{sagaḷyānnī vāhtukīcyā niymāñcaṁ pālan kelaṁ pāhije.} — Everyone must observe the traffic rules.
  \item \emph{Recall:} say the Marathi for a tendency, then the Marathi for social, then say \emph{pālan karṇe} again
\end{itemize}

\begin{tcolorbox}[breakable,colback=teal!4,colframe=teal!35!black,title={Before you move on}]
What does \mr{पालन} \mr{करणे} mean? (\textquotedblleft{}To comply with; to observe\textquotedblright{}.) Say it once more.
\end{tcolorbox}
\section[\texorpdfstring{\mr{उल्लंघन} \mr{करणे} (ullaṅghan karṇe)}{उल्लंघन करणे (ullaṅghan karṇe)}]{\mr{उल्लंघन} \mr{करणे} (ullaṅghan karṇe) — to violate; to break (a rule)}

\label{lesson:MR-C368-ullanghan-karne}



\begin{quote}
Before the new one: say the Marathi for social, then the Marathi for to comply with.
\end{quote}

\subsection*{You'll want to know: \mr{उल्लंघन} \mr{करणे}}
\textbf{\mr{उल्लंघन} \mr{करणे}} — \emph{ullaṅghan karṇe} — \textquotedblleft{}to violate; to break (a rule)\textquotedblright{}.

\begin{quote}
\mr{त्याने} \mr{नियमांचं} \mr{उल्लंघन} \mr{केलं}, \mr{म्हणून} \mr{त्याला} \mr{दंड} \mr{झाला}.
\end{quote}

\begin{quote}
\emph{tyāne niymāñcaṁ ullaṅghan kelaṁ, mhaṇūn tyālā daṇḍ jhālā.}
\end{quote}

\begin{quote}
He broke the rules, so he was fined.
\end{quote}

Past, with -\mr{ने} (he …): \mr{त्याने} \mr{उल्लंघन} \mr{केलं}, \emph{tyāne ullaṅghan kelaṁ}.

The rule takes -\mr{चं}: \mr{कायद्याचं} \mr{उल्लंघन}.

\begin{cousinweb}
\textbf{\mr{उल्लंघन} \mr{करणे}} (\emph{ullaṅghan karṇe}) joins the tatsama \mr{उल्लंघन}, Sanskrit \emph{ullaṅghana}, \textquotedblleft{}leaping over\textquotedblright{} (CDIAL 2365), to \mr{करणे}, \textquotedblleft{}to do\textquotedblright{}.
\end{cousinweb}

\begin{grammarlens}[title={What you've built}]
One more word for giving reasons and opinions.
\end{grammarlens}

\subsection*{Your turn}
\begin{itemize}
\raggedright
  \item \emph{Say it:} \emph{ullaṅghan karṇe}
  \item \emph{Say it:} \emph{ullaṅghan karṇe}, once more
  \item \emph{Say it:} \emph{tyāne niymāñcaṁ ullaṅghan kelaṁ, mhaṇūn tyālā daṇḍ jhālā.} — He broke the rules, so he was fined.
  \item \emph{Recall:} say the Marathi for social, then the Marathi for to comply with, then say \emph{ullaṅghan karṇe} again
\end{itemize}

\begin{tcolorbox}[breakable,colback=teal!4,colframe=teal!35!black,title={Before you move on}]
What does \mr{उल्लंघन} \mr{करणे} mean? (\textquotedblleft{}To violate; to break (a rule)\textquotedblright{}.) Say it once more.
\end{tcolorbox}
\section[\texorpdfstring{\mr{बजावणे} (bajāvṇe)}{बजावणे (bajāvṇe)}]{\mr{बजावणे} (bajāvṇe) — to discharge (a duty)}

\label{lesson:MR-C368-bajavne}



\begin{quote}
Before the new one: say the Marathi for to comply with, then the Marathi for to violate.
\end{quote}

\subsection*{You'll want to know: \mr{बजावणे}}
\textbf{\mr{बजावणे}} — \emph{bajāvṇe} — \textquotedblleft{}to discharge (a duty)\textquotedblright{}.

\begin{quote}
\mr{पोलिसाने} \mr{आपलं} \mr{कर्तव्य} \mr{प्रामाणिकपणे} \mr{बजावलं}.
\end{quote}

\begin{quote}
\emph{polisāne āplaṁ kartavya prāmāṇikpaṇe bajāvlaṁ.}
\end{quote}

\begin{quote}
The policeman did his duty honestly.
\end{quote}

Past, with -\mr{ने} (he …): \mr{त्याने} \mr{कर्तव्य} \mr{बजावलं}, \emph{tyāne kartavya bajāvlaṁ}.

\mr{कर्तव्य} \mr{बजावणे}, to do one's duty; also to warn sternly: \mr{त्याला} \mr{बजावलं}.

\begin{cousinweb}
\textbf{\mr{बजावणे}} (\emph{bajāvṇe}) is a Marathi verb whose origin is not settled.
\end{cousinweb}

\begin{grammarlens}[title={What you've built}]
One more word for giving reasons and opinions.
\end{grammarlens}

\subsection*{Your turn}
\begin{itemize}
\raggedright
  \item \emph{Say it:} \emph{bajāvṇe}
  \item \emph{Say it:} \emph{bajāvṇe}, once more
  \item \emph{Say it:} \emph{polisāne āplaṁ kartavya prāmāṇikpaṇe bajāvlaṁ.} — The policeman did his duty honestly.
  \item \emph{Recall:} say the Marathi for to comply with, then the Marathi for to violate, then say \emph{bajāvṇe} again
\end{itemize}

\begin{tcolorbox}[breakable,colback=teal!4,colframe=teal!35!black,title={Before you move on}]
What does \mr{बजावणे} mean? (\textquotedblleft{}To discharge (a duty)\textquotedblright{}.) Say it once more.
\end{tcolorbox}
\section[\texorpdfstring{\mr{सोपवणे} (sopavṇe)}{सोपवणे (sopavṇe)}]{\mr{सोपवणे} (sopavṇe) — to entrust; to hand over}

\label{lesson:MR-C368-sopavne}



\begin{quote}
Before the new one: say the Marathi for to violate, then the Marathi for to discharge.
\end{quote}

\subsection*{You'll want to know: \mr{सोपवणे}}
\textbf{\mr{सोपवणे}} — \emph{sopavṇe} — \textquotedblleft{}to entrust; to hand over\textquotedblright{}.

\begin{quote}
\mr{आईने} \mr{घराची} \mr{जबाबदारी} \mr{मोठ्या} \mr{मुलीवर} \mr{सोपवली}.
\end{quote}

\begin{quote}
\emph{āīne gharācī jabābdārī moṭhyā mulīvar sopavlī.}
\end{quote}

\begin{quote}
Mother entrusted the running of the house to her elder daughter.
\end{quote}

Past, with -\mr{ने} (he …): \mr{त्याने} \mr{सोपवलं}, \emph{tyāne sopavlaṁ}.

The person entrusted takes -\mr{वर} or -\mr{कडे}.

\begin{cousinweb}
\textbf{\mr{सोपवणे}} (\emph{sopavṇe}) is built on an older verb, Marathi \emph{sõpṇẽ}, \textquotedblleft{}hands over\textquotedblright{}, which CDIAL lists under Sanskrit \emph{samarpayati}, \textquotedblleft{}delivers over\textquotedblright{} (CDIAL 13192).
\end{cousinweb}

\begin{grammarlens}[title={What you've built}]
One more word for giving reasons and opinions.
\end{grammarlens}

\subsection*{Your turn}
\begin{itemize}
\raggedright
  \item \emph{Say it:} \emph{sopavṇe}
  \item \emph{Say it:} \emph{sopavṇe}, once more
  \item \emph{Say it:} \emph{āīne gharācī jabābdārī moṭhyā mulīvar sopavlī.} — Mother entrusted the running of the house to her elder daughter.
  \item \emph{Recall:} say the Marathi for to violate, then the Marathi for to discharge, then say \emph{sopavṇe} again
\end{itemize}

\begin{tcolorbox}[breakable,colback=teal!4,colframe=teal!35!black,title={Before you move on}]
What does \mr{सोपवणे} mean? (\textquotedblleft{}To entrust; to hand over\textquotedblright{}.) Say it once more.
\end{tcolorbox}
\section[\texorpdfstring{\mr{हाताळणे} (hātāḷṇe)}{हाताळणे (hātāḷṇe)}]{\mr{हाताळणे} (hātāḷṇe) — to handle}

\label{lesson:MR-C368-hatalne}



\begin{quote}
Before the new one: say the Marathi for to discharge, then the Marathi for to entrust.
\end{quote}

\subsection*{You'll want to know: \mr{हाताळणे}}
\textbf{\mr{हाताळणे}} — \emph{hātāḷṇe} — \textquotedblleft{}to handle\textquotedblright{}.

\begin{quote}
\mr{मीराने} \mr{तो} \mr{प्रश्न} \mr{शांतपणे} \mr{हाताळला}.
\end{quote}

\begin{quote}
\emph{mīrāne to praśna śāntpaṇe hātāḷlā.}
\end{quote}

\begin{quote}
Meera handled the issue calmly.
\end{quote}

Past, with -\mr{ने} (he …): \mr{त्याने} \mr{हाताळलं}, \emph{tyāne hātāḷlaṁ}.

Built on \mr{हात}, hand.

\begin{cousinweb}
\textbf{\mr{हाताळणे}} (\emph{hātāḷṇe}) is a Marathi verb built on \mr{हात}, \textquotedblleft{}hand\textquotedblright{}, which is inherited from Sanskrit \emph{hasta}, \textquotedblleft{}hand\textquotedblright{} (CDIAL 14024).
\end{cousinweb}

\begin{grammarlens}[title={What you've built}]
One more word for giving reasons and opinions.
\end{grammarlens}

\subsection*{Your turn}
\begin{itemize}
\raggedright
  \item \emph{Say it:} \emph{hātāḷṇe}
  \item \emph{Say it:} \emph{hātāḷṇe}, once more
  \item \emph{Say it:} \emph{mīrāne to praśna śāntpaṇe hātāḷlā.} — Meera handled the issue calmly.
  \item \emph{Recall:} say the Marathi for to discharge, then the Marathi for to entrust, then say \emph{hātāḷṇe} again
\end{itemize}

\begin{tcolorbox}[breakable,colback=teal!4,colframe=teal!35!black,title={Before you move on}]
What does \mr{हाताळणे} mean? (\textquotedblleft{}To handle\textquotedblright{}.) Say it once more.
\end{tcolorbox}
